\section{Jacobson rings}
\label{section-ring-jacobson}

\noindent
Let $R$ be a ring. The closed points of $\Spec(R)$ are the
maximal ideals of $R$. Often rings which occur naturally in algebraic
geometry have lots of maximal ideals. For example finite type algebras
over a field or over $\mathbf{Z}$. We will show that these
are examples of Jacobson rings.

\begin{definition}
\label{definition-ring-jacobson}
Let $R$ be a ring. We say that $R$ is a
{\it Jacobson ring} if every radical
ideal $I$ is the intersection of the
maximal ideals containing it.
\end{definition}

\begin{lemma}
\label{lemma-finite-type-field-Jacobson}
Any algebra of finite type over a field is Jacobson.
\end{lemma}

\begin{proof}
This follows from Theorem \ref{theorem-nullstellensatz}
and Definition \ref{definition-ring-jacobson}.
\end{proof}

\begin{lemma}
\label{lemma-jacobson-prime}
Let $R$ be a ring. If every prime ideal of $R$ is the
intersection of the maximal ideals containing it,
then $R$ is Jacobson.
\end{lemma}

\begin{proof}
This is immediately clear from the fact that
every radical ideal $I \subset R$ is the
intersection of the primes containing it.
See Lemma \ref{lemma-Zariski-topology}.
\end{proof}

\begin{lemma}
\label{lemma-jacobson}
A ring $R$ is Jacobson if and only if $\Spec(R)$
is Jacobson, see Topology,
Definition \ref{topology-definition-space-jacobson}.
\end{lemma}

\begin{proof}
Suppose $R$ is Jacobson. Let $Z \subset \Spec(R)$
be a closed subset. We have to show that the set of closed
points in $Z$ is dense in $Z$. Let $U \subset \Spec(R)$
be an open such that $U \cap Z$ is nonempty.
We have to show $Z \cap U$ contains a closed point
of $\Spec(R)$. We may
assume $U = D(f)$ as standard opens form a basis for the
topology on $\Spec(R)$. According to
Lemma \ref{lemma-Zariski-topology} we may assume that
$Z = V(I)$, where $I$ is a radical ideal. We see also
that $f \not \in I$. By assumption, there exists a
maximal ideal $\mathfrak m \subset R$ such that
$I \subset \mathfrak m$ but $f \not\in \mathfrak m$.
Hence $\mathfrak m \in D(f) \cap V(I) = U \cap Z$ as desired.

\medskip\noindent
Conversely, suppose that $\Spec(R)$ is Jacobson.
Let $I \subset R$ be a radical ideal. Let
$J = \cap_{I \subset \mathfrak m} \mathfrak m$
be the intersection of the maximal ideals containing $I$.
Clearly $J$ is a radical ideal, $V(J) \subset V(I)$, and
$V(J)$ is the smallest closed subset of $V(I)$ containing
all the closed points of $V(I)$. By assumption we see that
$V(J) = V(I)$. But Lemma \ref{lemma-Zariski-topology}
shows there is a bijection between Zariski closed
sets and radical ideals, hence $I = J$ as desired.
\end{proof}

\begin{lemma}
\label{lemma-characterize-jacobson}
Let $R$ be a ring. If $R$ is not Jacobson there exist
a prime $\mathfrak p \subset R$, an element $f \in R$
such that the following hold
\begin{enumerate}
\item $\mathfrak p$ is not a maximal ideal,
\item $f \not \in \mathfrak p$,
\item $V(\mathfrak p) \cap D(f) = \{\mathfrak p\}$, and
\item $(R/\mathfrak p)_f$ is a field.
\end{enumerate}
On the other hand, if $R$ is Jacobson, then for any pair $(\mathfrak p, f)$
such that (1) and (2) hold the set $V(\mathfrak p) \cap D(f)$ is
infinite.
\end{lemma}

\begin{proof}
Assume $R$ is not Jacobson.
By Lemma \ref{lemma-jacobson} this means there exists an
closed subset $T \subset \Spec(R)$
whose set $T_0 \subset T$ of closed points is not dense in $T$.
Choose an $f \in R$ such that $T_0 \subset V(f)$ but
$T \not \subset V(f)$. Note that $T \cap D(f)$
is homeomorphic to $\Spec((R/I)_f)$ if $T = V(I)$, see
Lemmas \ref{lemma-spec-closed} and \ref{lemma-standard-open}.
As any ring has a maximal ideal
(Lemma \ref{lemma-Zariski-topology}) we can choose a closed point $t$ of
space $T \cap D(f)$. Then $t$ corresponds to a prime ideal
$\mathfrak p \subset R$ which is not maximal (as $t \not \in T_0$).
Thus (1) holds. By construction $f \not \in \mathfrak p$, hence (2).
As $t$ is a closed point of $T \cap D(f)$ we see that
$V(\mathfrak p) \cap D(f) = \{\mathfrak p\}$, i.e., (3) holds. Hence we
conclude that $(R/\mathfrak p)_f$ is a domain whose
spectrum has one point, hence (4) holds
(for example combine Lemmas \ref{lemma-characterize-local-ring} and
\ref{lemma-minimal-prime-reduced-ring}).

\medskip\noindent
Conversely, suppose that $R$ is Jacobson and $(\mathfrak p, f)$
satisfy (1) and (2). If
$V(\mathfrak p) \cap D(f) =
\{\mathfrak p, \mathfrak q_1, \ldots, \mathfrak q_t\}$
then $\mathfrak p \not = \mathfrak q_i$
implies there exists an element $g \in R$ such that $g \not \in \mathfrak p$
but $g \in \mathfrak q_i$ for all $i$. Hence
$V(\mathfrak p) \cap D(fg) = \{\mathfrak p\}$ which
is impossible since each locally closed subset of $\Spec(R)$
contains at least one closed point as $\Spec(R)$ is
a Jacobson topological space.
\end{proof}

\begin{lemma}
\label{lemma-pid-jacobson}
The ring $\mathbf{Z}$ is a Jacobson ring.
More generally, let $R$ be a ring such that
\begin{enumerate}
\item $R$ is a domain,
\item $R$ is Noetherian,
\item any nonzero prime ideal is a maximal ideal, and
\item $R$ has infinitely many maximal ideals.
\end{enumerate}
Then $R$ is a Jacobson ring.
\end{lemma}

\begin{proof}
Let $R$ satisfy (1), (2), (3) and (4). The statement
means that $(0) = \bigcap_{\mathfrak m \subset R} \mathfrak m$.
Since $R$ has infinitely many maximal ideals it suffices to
show that any nonzero $x \in R$ is contained in at most
finitely many maximal ideals, in other words that $V(x)$ is finite.
By Lemma \ref{lemma-spec-closed}
we see that $V(x)$ is homeomorphic to $\Spec(R/xR)$.
By assumption (3) every prime of $R/xR$ is minimal and hence
corresponds to an irreducible component of $\Spec(R/xR)$
(Lemma \ref{lemma-irreducible}).
As $R/xR$ is Noetherian, the topological space $\Spec(R/xR)$
is Noetherian (Lemma \ref{lemma-Noetherian-topology})
and has finitely many irreducible components
(Topology, Lemma \ref{topology-lemma-Noetherian}).
Thus $V(x)$ is finite as desired.
\end{proof}

\begin{example}
\label{example-infinite-product-fields-jacobson}
Let $A$ be an infinite set.
For each $\alpha \in A$, let $k_\alpha$ be a field.
We claim that $R = \prod_{\alpha\in A} k_\alpha$ is Jacobson.
First, note that any element $f \in R$ has the form
$f = ue$, with $u \in R$ a unit and $e\in R$ an idempotent
(left to the reader).
Hence $D(f) = D(e)$, and $R_f = R_e = R/(1-e)$ is a quotient of $R$.
Actually, any ring with this property is Jacobson.
Namely, say $\mathfrak p \subset R$ is a prime ideal
and $f \in R$, $f \not \in \mathfrak p$. We have to find
a maximal ideal $\mathfrak m$ of $R$ such that
$\mathfrak p \subset \mathfrak m$ and $f \not\in \mathfrak m$.
Because $R_f$ is a quotient of $R$ we see that any maximal
ideal of $R_f$ corresponds to a maximal ideal of $R$
not containing $f$. Hence the result follows
by choosing a maximal ideal of $R_f$ containing $\mathfrak p R_f$.
\end{example}

\begin{example}
\label{example-not-jacobson}
A domain $R$ with finitely many maximal ideals
$\mathfrak m_i$, $i = 1, \ldots, n$ is not a
Jacobson ring, except when it is a field.
Namely, in this case $(0)$ is not the intersection
of the maximal ideals $(0) \not =
\mathfrak m_1 \cap \mathfrak m_2 \cap \ldots \cap \mathfrak m_n
\supset \mathfrak m_1 \cdot \mathfrak m_2 \cdot \ldots
\cdot \mathfrak m_n \not = 0$.
In particular a discrete valuation ring, or any local ring with
at least two prime ideals is not a Jacobson
ring.
\end{example}

\begin{lemma}
\label{lemma-finite-residue-extension-closed}
Let $R \to S$ be a ring map.
Let $\mathfrak m \subset R$ be a maximal ideal.
Let $\mathfrak q \subset S$ be a prime ideal
lying over $\mathfrak m$ such that $\kappa(\mathfrak q)/\kappa(\mathfrak m)$
is an algebraic field extension.
Then $\mathfrak q$ is a maximal ideal of $S$.
\end{lemma}

\begin{proof}
Consider the diagram
$$
\xymatrix{
S \ar[r] & S/\mathfrak q \ar[r] & \kappa(\mathfrak q) \\
R \ar[r] \ar[u] & R/\mathfrak m \ar[u]
}
$$
We see that $\kappa(\mathfrak m) \subset S/\mathfrak q \subset
\kappa(\mathfrak q)$. Because the field extension
$\kappa(\mathfrak m) \subset \kappa(\mathfrak q)$
is algebraic, any ring between $\kappa(\mathfrak m)$
and $\kappa(\mathfrak q)$ is a field
(Fields, Lemma \ref{fields-lemma-subalgebra-algebraic-extension-field}).
Thus $S/\mathfrak q$ is a field, and a posteriori equal
to $\kappa(\mathfrak q)$.
\end{proof}

\begin{lemma}
\label{lemma-dimension}
Suppose that $k$ is a field and suppose that $V$ is a nonzero vector
space over $k$. Assume the dimension of $V$ (which is a cardinal number)
is smaller than the cardinality of $k$. Then for any linear operator
$T : V \to V$ there exists some monic polynomial $P(t) \in k[t]$ such that
$P(T)$ is not invertible.
\end{lemma}

\begin{proof}
If not then $V$ inherits the structure of a vector space over
the field $k(t)$. But the dimension of $k(t)$ over $k$ is
at least the cardinality of $k$ for example due to the fact that the elements
$\frac{1}{t - \lambda}$ are $k$-linearly independent.
\end{proof}

\noindent
Here is another version of Hilbert's Nullstellensatz.

\begin{theorem}
\label{theorem-uncountable-nullstellensatz}
Let $k$ be a field. Let $S$ be a $k$-algebra generated over $k$
by the elements $\{x_i\}_{i \in I}$. Assume the cardinality of $I$
is smaller than the cardinality of $k$. Then
\begin{enumerate}
\item for all maximal ideals $\mathfrak m \subset S$ the field
extension $\kappa(\mathfrak m)/k$
is algebraic, and
\item $S$ is a Jacobson ring.
\end{enumerate}
\end{theorem}

\begin{proof}
If $I$ is finite then the result follows from the Hilbert Nullstellensatz,
Theorem \ref{theorem-nullstellensatz}. In the rest of the proof we assume
$I$ is infinite. It suffices to prove the result for
$\mathfrak m \subset k[\{x_i\}_{i \in I}]$ maximal in the polynomial
ring on variables $x_i$, since $S$ is a quotient of this.
As $I$ is infinite the set
of monomials $x_{i_1}^{e_1} \ldots x_{i_r}^{e_r}$, $i_1, \ldots, i_r \in I$
and $e_1, \ldots, e_r \geq 0$ has cardinality at most equal to the
cardinality of $I$. Because the cardinality of $I \times \ldots \times I$
is the cardinality of $I$, and also the cardinality of
$\bigcup_{n \geq 0} I^n$ has the same cardinality.
(If $I$ is finite, then this is not true and
in that case this proof only works if $k$ is uncountable.)

\medskip\noindent
To arrive at a contradiction pick $T \in \kappa(\mathfrak m)$ transcendental
over $k$. Note that the $k$-linear map
$T : \kappa(\mathfrak m) \to \kappa(\mathfrak m)$
given by multiplication by $T$ has the property that $P(T)$ is invertible
for all monic polynomials $P(t) \in k[t]$.
Also, $\kappa(\mathfrak m)$ has dimension at most the cardinality of $I$
over $k$ since it is a quotient of the vector space
$k[\{x_i\}_{i \in I}]$ over $k$ (whose dimension is $\# I$ as we saw above).
This is impossible by Lemma \ref{lemma-dimension}.

\medskip\noindent
To show that $S$ is Jacobson we argue as follows. If not then
there exists a prime $\mathfrak q \subset S$ and an element $f \in S$,
$f \not \in \mathfrak q$ such that $\mathfrak q$ is not maximal
and $(S/\mathfrak q)_f$ is a field, see
Lemma \ref{lemma-characterize-jacobson}.
But note that $(S/\mathfrak q)_f$ is generated by at most $\# I + 1$ elements.
Hence the field extension $(S/\mathfrak q)_f/k$ is algebraic
(by the first part of the proof).
This implies that $\kappa(\mathfrak q)$ is an algebraic extension of $k$
hence $\mathfrak q$ is maximal by
Lemma \ref{lemma-finite-residue-extension-closed}. This contradiction
finishes the proof.
\end{proof}

\begin{lemma}
\label{lemma-base-change-Jacobson}
Let $k$ be a field. Let $S$ be a $k$-algebra.
For any field extension $K/k$ whose cardinality is larger
than the cardinality of $S$ we have
\begin{enumerate}
\item for every maximal ideal $\mathfrak m$ of $S_K$ the field
$\kappa(\mathfrak m)$ is algebraic over $K$, and
\item $S_K$ is a Jacobson ring.
\end{enumerate}
\end{lemma}

\begin{proof}
Choose $k \subset K$ such that the cardinality of $K$ is greater
than the cardinality of $S$. Since the elements of $S$ generate
the $K$-algebra $S_K$ we see that
Theorem \ref{theorem-uncountable-nullstellensatz}
applies.
\end{proof}

\begin{example}
\label{example-countable-trick-does-not-work}
The trick in the proof of
Theorem \ref{theorem-uncountable-nullstellensatz}
really does not work if $k$ is a countable field and $I$ is countable too.
Let $k$ be a countable field. Let $x$ be a variable,
and let $k(x)$ be the field of rational functions in $x$.
Consider the polynomial algebra $R = k[x, \{x_f\}_{f \in k[x]-\{0\}}]$.
Let $I = (\{fx_f - 1\}_{f\in k[x] - \{0\}})$. Note that
$I$ is a proper ideal in $R$.
Choose a maximal ideal $I \subset \mathfrak m$.
Then $k \subset R/\mathfrak m$ is isomorphic to
$k(x)$, and is not algebraic over $k$.
\end{example}

\begin{lemma}
\label{lemma-Jacobson-invert-element}
Let $R$ be a Jacobson ring. Let $f \in R$. The ring $R_f$ is Jacobson and
maximal ideals of $R_f$ correspond to maximal ideals of $R$ not containing $f$.
\end{lemma}

\begin{proof}
By Topology, Lemma \ref{topology-lemma-jacobson-inherited}
we see that $D(f) = \Spec(R_f)$ is Jacobson and
that closed points of $D(f)$
correspond to closed points in $\Spec(R)$
which happen to lie in $D(f)$. Thus $R_f$ is Jacobson by
Lemma \ref{lemma-jacobson}.
\end{proof}

\begin{example}
\label{example-localize-not-preserve-closed-points}
Here is a simple example that shows Lemma \ref{lemma-Jacobson-invert-element}
to be false if $R$ is not Jacobson.
Consider the ring $R = \mathbf{Z}_{(2)}$, i.e., the localization
of $\mathbf{Z}$ at the prime $(2)$. The localization of $R$ at
the element $2$ is isomorphic to $\mathbf{Q}$, in a formula:
$R_2 \cong \mathbf{Q}$. Clearly the map $R \to R_2$ maps the
closed point of $\Spec(\mathbf{Q})$ to the generic point
of $\Spec(R)$.
\end{example}

\begin{example}
\label{example-infinite-localize-not-preserve-closed-points}
Here is a simple example that shows
Lemma \ref{lemma-Jacobson-invert-element}
is false if $R$ is Jacobson but we localize at infinitely
many elements.
Namely, let $R = \mathbf{Z}$ and consider the localization
$(R \setminus \{0\})^{-1}R \cong \mathbf{Q}$
of $R$ at the set of all nonzero elements.
Clearly the map $\mathbf{Z} \to \mathbf{Q}$ maps the
closed point of $\Spec(\mathbf{Q})$ to the generic point
of $\Spec(\mathbf{Z})$.
\end{example}

\begin{lemma}
\label{lemma-Jacobson-mod-ideal}
Let $R$ be a Jacobson ring. Let $I \subset R$ be an ideal.
The ring $R/I$ is Jacobson and maximal ideals
of $R/I$ correspond to maximal ideals of $R$ containing $I$.
\end{lemma}

\begin{proof}
The proof is the same as the proof of
Lemma \ref{lemma-Jacobson-invert-element}.
\end{proof}

\begin{lemma}
\label{lemma-silly-jacobson}
Let $R$ be a Jacobson ring. Let $K$ be a field. Let $R \subset K$ and
$K$ is of finite type over $R$. Then $R$ is a field and $K/R$
is a finite field extension.
\end{lemma}

\begin{proof}
First note that $R$ is a domain.
By Lemma \ref{lemma-field-finite-type-over-domain}
we see that $R_f$ is a field and $K/R_f$ is a finite field extension
for some nonzero $f \in R$. Hence $(0)$ is a maximal ideal of $R_f$
and by
Lemma \ref{lemma-Jacobson-invert-element}
we conclude $(0)$ is a maximal ideal of $R$.
\end{proof}

\begin{proposition}
\label{proposition-Jacobson-permanence}
Let $R$ be a Jacobson ring. Let $R \to S$ be a
ring map of finite type. Then
\begin{enumerate}
\item The ring $S$ is Jacobson.
\item The map $\Spec(S) \to \Spec(R)$ transforms
closed points to closed points.
\item For $\mathfrak m' \subset S$ maximal lying over $\mathfrak m \subset R$
the field extension $\kappa(\mathfrak m')/\kappa(\mathfrak m)$
is finite.
\end{enumerate}
\end{proposition}

\begin{proof}
Let $\mathfrak m' \subset S$ be a maximal ideal and
$R \cap \mathfrak m' = \mathfrak m$.
Then $R/\mathfrak m \to S/\mathfrak m'$ satisfies
the conditions of Lemma \ref{lemma-silly-jacobson}
by Lemma \ref{lemma-Jacobson-mod-ideal}.
Hence $R/\mathfrak m$ is a field and
$\mathfrak m$ a maximal ideal and the induced
residue field extension is finite. This proves (2) and (3).  

\medskip\noindent
If $S$ is not Jacobson, then by Lemma \ref{lemma-characterize-jacobson} there
exists a non-maximal prime ideal $\mathfrak q$ of $S$ and an
$g \in S$, $g \not\in \mathfrak q$ such that $(S/\mathfrak q)_g$ is a field.
To arrive at a contradiction we show that $\mathfrak q$ is a maximal ideal.
Let $\mathfrak p = \mathfrak q \cap R$. Then
$R/\mathfrak p \to (S/\mathfrak q)_g$ satisfies the conditions of
Lemma \ref{lemma-silly-jacobson} by
Lemma \ref{lemma-Jacobson-mod-ideal}.
Hence $R/\mathfrak p$ is a field and the field extension
$\kappa(\mathfrak p) \to (S/\mathfrak q)_g = \kappa(\mathfrak q)$ is
finite, thus algebraic. Then $\mathfrak q$ is a maximal ideal of $S$ by
Lemma \ref{lemma-finite-residue-extension-closed}. Contradiction.
\end{proof}

\begin{lemma}
\label{lemma-corollary-jacobson}
Any finite type algebra over $\mathbf{Z}$ is Jacobson.
\end{lemma}

\begin{proof}
Combine Lemma \ref{lemma-pid-jacobson} and
Proposition \ref{proposition-Jacobson-permanence}.
\end{proof}

\begin{lemma}
\label{lemma-image-finite-type-map-Jacobson-rings}
Let $R \to S$ be a finite type ring map of Jacobson rings.
Denote $X = \Spec(R)$ and $Y = \Spec(S)$.
Write $f : Y \to X$ the induced
map of spectra. Let $E \subset Y = \Spec(S)$ be a
constructible set. Denote with a subscript ${}_0$ the set
of closed points of a topological space.
\begin{enumerate}
\item We have $f(E)_0 = f(E_0) = X_0 \cap f(E)$.
\item A point $\xi \in X$ is in $f(E)$ if and only if
$\overline{\{\xi\}} \cap f(E_0)$ is dense in $\overline{\{\xi\}}$.
\end{enumerate}
\end{lemma}

\begin{proof}
We have a commutative diagram of continuous maps
$$
\xymatrix{
E \ar[r] \ar[d] & Y \ar[d] \\
f(E) \ar[r] & X
}
$$
Suppose $x \in f(E)$ is closed in $f(E)$. Then $f^{-1}(\{x\})\cap E$
is nonempty and closed in $E$. Applying
Topology, Lemma \ref{topology-lemma-jacobson-inherited}
to both inclusions
$$
f^{-1}(\{x\}) \cap E \subset E \subset Y
$$
we find there exists a point $y \in f^{-1}(\{x\}) \cap E$ which is
closed in $Y$. In other words, there exists $y \in Y_0$ and $y \in E_0$
mapping to $x$. Hence $x \in f(E_0)$.
This proves that $f(E)_0 \subset f(E_0)$.
Proposition \ref{proposition-Jacobson-permanence} implies that
$f(E_0) \subset X_0 \cap f(E)$. The inclusion
$X_0 \cap f(E) \subset f(E)_0$ is trivial. This proves the
first assertion.

\medskip\noindent
Suppose that $\xi \in f(E)$. According to
Lemma \ref{lemma-characterize-image-finite-type}
the set $f(E) \cap \overline{\{\xi\}}$ contains a dense
open subset of $\overline{\{\xi\}}$. Since $X$ is Jacobson
we conclude that $f(E) \cap \overline{\{\xi\}}$ contains a
dense set of closed points, see Topology,
Lemma \ref{topology-lemma-jacobson-inherited}.
We conclude by part (1) of the lemma.

\medskip\noindent
On the other hand, suppose that $\overline{\{\xi\}} \cap f(E_0)$
is dense in $\overline{\{\xi\}}$. By
Lemma \ref{lemma-constructible-is-image}
there exists a ring map $S \to S'$ of finite presentation
such that $E$ is the image of $Y' := \Spec(S') \to Y$.
Then $E_0$ is the image of $Y'_0$ by the first part of the
lemma applied to the ring map $S \to S'$. Thus we may assume that
$E = Y$ by replacing $S$ by $S'$. Suppose $\xi$ corresponds
to $\mathfrak p \subset R$. Consider the diagram
$$
\xymatrix{
S \ar[r] & S/\mathfrak p S \\
R \ar[r] \ar[u] & R/\mathfrak p \ar[u]
}
$$
This diagram and the density of $f(Y_0) \cap V(\mathfrak p)$
in $V(\mathfrak p)$
shows that the morphism $R/\mathfrak p \to S/\mathfrak p S$
satisfies condition (2) of
Lemma \ref{lemma-domain-image-dense-set-points-generic-point}.
Hence we conclude
there exists a prime $\overline{\mathfrak q} \subset S/\mathfrak pS$
mapping to $(0)$. In other words the inverse image $\mathfrak q$
of $\overline{\mathfrak q}$ in $S$ maps to $\mathfrak p$ as desired.
\end{proof}

\noindent
The conclusion of the lemma above is that we can read off
the image of $f$ from the set of closed points of the image.
This is a little nicer in case the map is of finite presentation
because then we know that images of a constructible is constructible.
Before we state it we introduce some notation.
Denote $\text{Constr}(X)$ the set of constructible sets.
Let $R \to S$ be a ring map.
Denote $X = \Spec(R)$ and $Y = \Spec(S)$.
Write $f : Y \to X$ the induced map of spectra.
Denote with a subscript ${}_0$ the set
of closed points of a topological space.


\begin{lemma}
\label{lemma-conclude-jacobson-Noetherian}
With notation as above. Assume that $R$ is a Noetherian Jacobson ring.
Further assume $R \to S$ is of finite type.
There is a commutative diagram
$$
\xymatrix{
\text{Constr}(Y) \ar[r]^{E \mapsto E_0} \ar[d]^{E \mapsto f(E)} &
\text{Constr}(Y_0) \ar[d]^{E \mapsto f(E)} \\
\text{Constr}(X) \ar[r]^{E \mapsto E_0} &
\text{Constr}(X_0)
}
$$
where the horizontal arrows are the bijections from
Topology, Lemma \ref{topology-lemma-jacobson-equivalent-constructible}.
\end{lemma}

\begin{proof}
Since $R \to S$ is of finite type, it is of finite presentation,
see Lemma \ref{lemma-Noetherian-finite-type-is-finite-presentation}.
Thus the image of a constructible set in $X$ is constructible
in $Y$ by Chevalley's theorem
(Theorem \ref{theorem-chevalley}).
Combined with
Lemma \ref{lemma-image-finite-type-map-Jacobson-rings}
the lemma follows.
\end{proof}

\noindent
To illustrate the use of Jacobson rings, we give the following two examples.

\begin{example}
\label{example-product-matrices-zero}
Let $k$ be a field. The space $\Spec(k[x, y]/(xy))$
has two irreducible components: namely the $x$-axis and the
$y$-axis. As a generalization, let
$$
R = k[x_{11}, x_{12}, x_{21}, x_{22}, y_{11}, y_{12}, y_{21}, y_{22}]/
\mathfrak a,
$$
where $\mathfrak a$ is the ideal in
$k[x_{11}, x_{12}, x_{21}, x_{22}, y_{11}, y_{12}, y_{21}, y_{22}]$
generated by the entries of the $2 \times 2$ product matrix
$$
\left(
\begin{matrix}
x_{11} & x_{12}\\
x_{21} & x_{22}
\end{matrix}
\right)
 \left(
\begin{matrix}
y_{11} & y_{12}\\
y_{21} & y_{22}
\end{matrix}
\right).
$$
In this example we will describe $\Spec(R)$.

\medskip\noindent
To prove the statement about $\Spec(k[x, y]/(xy))$ we argue as follows.
If $\mathfrak p \subset k[x, y]$ is any ideal containing $xy$, then either
$x$ or $y$ would be contained in $\mathfrak p$. Hence the minimal such
prime ideals are just $(x)$ and $(y)$. In case $k$ is
algebraically closed, the $\text{max-Spec}$ of these components
can then be visualized as the point sets of $y$- and $x$-axis.

\medskip\noindent
For the generalization, note that we may identify the closed
points of the spectrum of
$k[x_{11}, x_{12}, x_{21}, x_{22}, y_{11}, y_{12}, y_{21}, y_{22}])$
with the space of matrices
$$
\left\{ (X, Y) \in \text{Mat}(2, k)\times \text{Mat}(2, k) \mid
X = \left(
\begin{matrix}
x_{11} & x_{12}\\
x_{21} & x_{22}
\end{matrix}
\right),
Y= \left(
\begin{matrix}
y_{11} & y_{12}\\
y_{21} & y_{22}
\end{matrix}
\right)
\right\}
$$
at least if $k$ is algebraically closed.
Now define a group action of
$\text{GL}(2, k)\times \text{GL}(2, k)\times \text{GL}(2, k)$
on the space of matrices $\{(X, Y)\}$ by
$$
(g_1, g_2, g_3) \times (X, Y) \mapsto ((g_1Xg_2^{-1}, g_2Yg_3^{-1})).
$$
Here, also observe that the algebraic set
$$
\text{GL}(2, k)\times \text{GL}(2, k)\times \text{GL}(2, k) \subset
\text{Mat}(2, k)\times \text{Mat}(2, k) \times \text{Mat}(2, k)
$$
is irreducible since it is the max spectrum of the domain
$$
k[x_{11}, x_{12}, \ldots, z_{21}, z_{22}, (x_{11}x_{22}-x_{12}x_{21})^{-1}
, (y_{11}y_{22}-y_{12}y_{21})^{-1}, (z_{11}z_{22}-z_{12}z_{21})^{-1}].
$$
Since the image of irreducible an algebraic set is still
irreducible, it suffices to classify the orbits of the set
$\{(X, Y)\in \text{Mat}(2, k)\times \text{Mat}(2, k)|XY = 0\}$ and take their
closures. From standard linear algebra, we are reduced to the
following three cases:
\begin{enumerate}
\item $\exists (g_1, g_2)$ such that $g_1Xg_2^{-1} = I_{2\times 2}$.
Then $Y$ is necessarily $0$, which as an algebraic set is
invariant under the group action. It follows that this orbit is
contained in the irreducible algebraic set defined by the prime
ideal $(y_{11}, y_{12}, y_{21}, y_{22})$. Taking the closure, we see
that $(y_{11}, y_{12}, y_{21}, y_{22})$ is actually a component.
\item $\exists (g_1, g_2)$ such that
$$
g_1Xg_2^{-1} = \left(
\begin{matrix}
1 & 0 \\
0 & 0
\end{matrix}
\right).
$$
This case occurs if and only if $X$ is a rank 1 matrix,
and furthermore, $Y$ is killed by such an $X$ if and only if
$$
x_{11}y_{11}+x_{12}y_{21} = 0; \quad x_{11}y_{12}+x_{12}y_{22} = 0;
$$
$$
x_{21}y_{11}+x_{22}y_{21} = 0; \quad x_{21}y_{12}+x_{22}y_{22} = 0.
$$
Fix a rank 1 $X$, such non zero $Y$'s satisfying the above
equations form an irreducible algebraic set for the following
reason($Y = 0$ is contained the previous case):
$0 = g_1Xg_2^{-1}g_2Y$ implies that
$$
g_2Y = \left(
\begin{matrix}
0 & 0 \\
y_{21}' & y_{22}'
\end{matrix}
\right).
$$
With a further $\text{GL}(2, k)$-action on the right by $g_3$,
$g_2Y$ can be brought into
$$
g_2Yg_3^{-1} = \left(
\begin{matrix}
0 & 0 \\
0 & 1
\end{matrix}
\right),
$$
and thus such $Y$'s form an irreducible algebraic set
isomorphic to the image of $\text{GL}(2, k)$ under this action. Finally,
notice that the ``rank 1" condition for $X$'s forms an open dense
subset of the irreducible algebraic set
$\det X = x_{11}x_{22} - x_{12}x_{21} = 0$.
It now follows that all the five equations define an irreducible component
$(x_{11}y_{11}+x_{12}y_{21}, x_{11}y_{12}+x_{12}y_{22}, x_{21}y_{11}
+x_{22}y_{21}, x_{21}y_{12}+x_{22}y_{22}, x_{11}x_{22}-x_{12}x_{21})$
in the open subset of the space of pairs of nonzero matrices.
It can be shown that the pair of equations
$\det X = 0$, $\det Y = 0$ cuts $\Spec(R)$ in an irreducible component
with the above locus an open dense subset.
\item $\exists (g_1, g_2)$ such that $g_1Xg_2^{-1} = 0$, or
equivalently, $X = 0$. Then $Y$ can be arbitrary and this component
is thus defined by $(x_{11}, x_{12}, x_{21}, x_{22})$.
\end{enumerate}
\end{example}


\begin{example}
\label{example-idempotent-matrices}
For another example, consider
$R = k[\{t_{ij}\}_{i, j = 1}^{n}]/\mathfrak a$, where $\mathfrak a$ is
the ideal generated by the entries of the product matrix $T^2-T$,
$T = (t_{ij})$. From linear algebra, we know that under the
$GL(n, k)$-action defined by $g, T \mapsto gTg^{-1}$, $T$ is
classified by the its rank and each $T$ is conjugate to some
$\text{diag}(1, \ldots, 1, 0, \ldots, 0)$, which has $r$ 1's and $n-r$ 0's.
Thus each orbit of such a $\text{diag}(1, \ldots, 1, 0, \ldots, 0)$ under the
group action forms an irreducible component and every idempotent
matrix is contained in one such orbit. Next we will show that any
two different orbits are necessarily disjoint. For this purpose we
only need to cook up polynomial functions that take different
values on different orbits. In characteristic 0 cases, such a
function can be taken to be
$f(t_{ij}) = trace(T) = \sum_{i = 1}^nt_{ii}$. In positive
characteristic cases, things are slightly more tricky since we
might have $trace(T) = 0$ even if $T \neq 0$. For instance, $char = 3$
$$
trace\left(
\begin{matrix}
1 & & \\
& 1 & \\
& & 1
\end{matrix}
\right) = 3 = 0
$$
Anyway, these components can be separated using other functions.
For instance, in the characteristic 3 case, $tr(\wedge^3T)$ takes
value 1 on the components corresponding to $diag(1, 1, 1)$ and 0 on
other components.
\end{example}





































